\chapter{Constitutive Relationships and Reduction}

The constitutive law itself is unchanged on the current PMG path:
\begin{itemize}
\item the same Mohr-Coulomb return mapping,
\item the same Davis reduction,
\item the same stress tensor \(S\),
\item the same consistent tangent \(DS\).
\end{itemize}

What changes on the mainline path is how those quantities are owned, cached,
and exposed to distributed PETSc operators.

\section{Mainline Constitutive Cycle}

For the committed benchmark, one constitutive evaluation still means:
\begin{enumerate}
\item reduce strength parameters for the current \(\lambda\),
\item compute strain from the current displacement field,
\item evaluate stresses and tangent blocks at integration points,
\item assemble the internal force,
\item build or refresh the tangent and regularized operators needed by Newton.
\end{enumerate}

That is the same cycle MATLAB authors already know.

\section{What The Mainline Builder Owns}

The mainline \code{ConstitutiveOperator} in
\petscmod{constitutive/problem.py} owns:
\begin{itemize}
\item the material arrays \code{c0}, \code{phi}, \code{psi}, \code{shear},
\code{bulk}, \code{lame},
\item the free-DOF mask \code{q\_mask},
\item the current stress and tangent state,
\item the owned tangent metadata supplied by
\code{OwnedTangentPattern},
\item cached PETSc matrices for tangent and regularized operators on that owned
pattern.
\end{itemize}

The widened class boundary is a PETSc requirement, not a constitutive change.

\section{Mainline Ownership Mode: Overlap}

For the committed 3D benchmark, the important constitutive fact is simple:
\begin{itemize}
\item \code{constitutive\_mode = overlap},
\item the owned tangent pattern is prepared with
\code{include\_unique = false},
\item each rank evaluates constitutive data directly on the overlap integration
points it needs for its owned rows,
\item no unique-owner gather or exchange stage participates in the mainline
run.
\end{itemize}

So the current benchmark should be read as an \emph{overlap-evaluated
constitutive path}, not as a unique-owner recovery path.

\section{Which Builder Interfaces Are Mainline}

Because the linear solver backend is \code{pmg\_shell} and does not prefer the
full-system operator, the mainline Newton path mainly consumes:
\begin{itemize}
\item \code{build\_F\_free(...)} or equivalent free-output residual builders,
\item \code{build\_F\_K\_tangent\_reduced\_free(...)} and related free-output
tangent builders,
\item \code{build\_K\_regularized(...)} for the regularized owned-row PETSc
matrix.
\end{itemize}

So although the class exposes many interfaces, the default 3D PMG route is the
free-output interface family backed by owned-row caches.

\section{What Stayed The Same For MATLAB Authors}

The most important parity statement is still:
\begin{itemize}
\item \code{reduction(lambda)} still means strength reduction,
\item strain still comes from the current displacement field,
\item \code{constitutive\_problem\_3D(...)} still computes the same stresses and
tangent blocks,
\item the internal force and tangent still represent the same nonlinear
mechanics.
\end{itemize}

The PETSc scaffolding around those steps exists because MPI ownership and solver
reuse are now explicit.

\section{Where To Edit}

\begin{itemize}
\item \petscmod{constitutive/reduction.py} for Davis reduction rules,
\item \petscmod{constitutive/problem.py} for stress, tangent, residual, and
regularized-operator exposure,
\item \petscmod{fem/distributed\_tangent.py} when the owned-row constitutive
interface changes.
\end{itemize}

\section{Where The Other Constitutive Modes Went}

The global and unique-owner constitutive modes, along with the BDDC-specific
builder interfaces, are moved to Appendix~\ref{app:constitutive}.
